\section{Zusammenfassung und Ausblick}

\begin{frame}{Zusammenfassung}
Wir haben kennengelernt:
\begin{itemize}
\item verschiedene \alert{Datentypen} (tw. \glqq High Level\grqq)
\item die wichtigsten \alert{Statements}
\item \alert{Funktion}sdeklaration und -Benutzung
\item \alert{Module} und Pakete
\item Fehler und \alert{Ausnahmen}, Behandlung selbiger
\item \alert{objektorientierte Programmierung}
\item einige h"aufig verwendete Standardmodule
\end{itemize}
\end{frame}

\begin{frame}{Offene Punkte}
Nicht behandelte, tw. fortgeschrittene Themen:
\begin{itemize}
\item Closures, Dekoratoren (Funktionswrapper)
\item Metaklassen
\item Weitere Standardmodule: Mail, WWW, XML, Zeit\&Datum, \dots $\rightarrow$~\href{http://docs.python.org/lib/modindex.html}{http://docs.python.org/lib/modindex.html}
\item Profiling, Debugging, Unittesting
\item Extending und Embedding: Python \& C/C++ $\rightarrow$~\href{http://docs.python.org/ext/ext.html}{http://docs.python.org/ext/ext.html}
\item Third Party-Module: Grafik, Webprogrammierung, Datenbanken, \dots $\rightarrow$~\href{http://pypi.python.org/pypi}{http://pypi.python.org/pypi}
\end{itemize}
\end{frame}

\begin{frame}{Web-Programmierung}
\begin{itemize}
\item CGI-Scripte: Modul \texttt{cgi} aus  Standardbibliothek
\item Webframeworks: Django, TurboGears, Pylons, \dots
\item Templatesysteme: Cheetah, Genshi, Jinja, \dots
\item Content Management Systeme (CMS): Zope, Plone, Skeletonz, \dots
\item Wikis: MoinMoin, \dots
\end{itemize}

\includegraphics[height=4cm]{images/moinmoin.png}
\end{frame}

\begin{frame}[fragile]{NumPy + SciPy + Matplotlib = Pylab}
Ein Ersatz f"ur MatLab: Matritzenrechnung, numerische Funktionen, Plotten, ...

\includegraphics[height=4cm]{images/matplotlib.png}
\includegraphics[height=4cm]{images/surface.jpg}

\begin{lstlisting}[style=Python, basicstyle=\small]
A = matrix([[1, 2], [2, 1]]); b = array([1, -1])
matshow(A)
(eigvals, eigvecs) = eig(A)
x = linalg.solve(A, b)
\end{lstlisting}
\end{frame}

\begin{frame}[fragile]{Und mehr...}
\begin{itemize}
\item \alert{ipython}: Eine komfortablere Python-Shell
\item \alert{Scons}: Ein Build-System "ahnlich Make/autoconf/automake oder Ant
\begin{itemize}
\item basiert auf Python, nutzt Python-Syntax
\item Built-in-support f"ur viele Programmiersprachen, LaTeX, \dots
\item plattformunabh"angig
\end{itemize}
\item Python und andere Programmiersprachen:
\begin{itemize}
\item Jython: Python-Code in der Java VM ausf"uhren
\item Ctypes: C-Libraries mit Python ansprechen (ab 2.5 in der stdlib)
\item SWIG: C- und C++ -Libraries mit Python ansprechen
\end{itemize}
\item \alert{PIL}: Python Imaging Library f"ur Bildmanipulation
\end{itemize}
\end{frame}


\begin{frame}{PyCologne}
\begin{columns}
\begin{column}{0.4\textwidth}
  \includegraphics[height=5cm]{images/pycologne7.pdf}
\end{column}
\begin{column}{0.6\textwidth}
  \alert{PyCologne}: Python User Group K\"oln
  \begin{itemize}
    \item Trifft sich jeden zweiten Mittwoch im Monat am Rechenzentrum der Uni K\"oln
    \item URL: \href{http://wiki.python.de/pyCologne}{http://wiki.python.de/pyCologne}
  \end{itemize}
\end{column}
\end{columns}
  
\end{frame}
